% Part: incompleteness
% Chapter: representability-in-q
% Section: representing-relations

\documentclass[../../../include/open-logic-section]{subfiles}

\begin{document}

\olfileid[pt]{inc}{req}{rel}

\olsection{Representação de relações}

Digamos o que significa uma \emph{relação} ser representável.

\begin{defn}
\ollabel{defn:representing-relations} Uma relação $R(x_0,\dots,x_k)$
sobre os números naturais é {\em representável em $\Th{Q}$} se existe
uma fórmula $!A_R(x_0,\dots,x_k)$ tal que, sempre que
$R(n_0,\dots,n_k)$ é verdadeira, $\Th{Q}$ prova
$!A_R(\num{n_0},\dots,\num{n_k})$ e, sempre que
$R(n_0,\dots,n_k)$ é falsa, $\Th{Q}$ prova $\lnot !A_R(\num{n_0},
\dots, \num{n_k})$.
\end{defn}

\begin{thm}
\ollabel{thm:representing-rels} Uma relação é representável em
$\Th{Q}$ se, e somente se, é computável.
\end{thm}

\begin{proof}
Para a direção direta, suponha que $R(x_0,\dots,x_k)$ seja
representada pela fórmula $!A_R(x_0,\dots,x_k)$. Eis um algoritmo para
computar $R$: com entrada $n_0$, \dots,~$n_k$, procure simultaneamente
uma prova de $!A_R(\num{n_0}, \dots, \num{n_k})$ e uma prova de
$\lnot !A_R(\num{n_0}, \dots, \num{n_k})$. Por hipótese, a busca
necessariamente encontrará uma ou outra; se encontrar a primeira,
responda ``sim'' e, caso contrário, responda ``não''.

Na outra direção, suponha que $R(x_0, \dots, x_k)$ seja computável.
Por definição, isso significa que a função $\Char{R}(x_0, \dots, x_k)$
é computável. Por \olref[int]{thm:representable-iff-comp}, $\Char{R}$
é representada por uma fórmula, digamos,
$!A_{\Char{R}}(x_0, \dots, x_k, y)$. Seja $!A_R(x_0, \dots, x_k)$ a
fórmula $!A_{\Char{R}}(x_0, \dots, x_k, \num{1})$. Então, para
quaisquer $n_0$, \dots,~$n_k$, se $R(n_0, \dots, n_k)$ é verdadeira,
então $\Char{R}(n_0, \dots, n_k) = 1$; nesse caso, $\Th{Q}$ prova
$!A_{\Char{R}}(\num{n_0}, \dots, \num{n_k}, \num{1})$ e, portanto,
$\Th{Q}$ prova $!A_R(\num{n_0}, \dots, \num{n_k})$. Por outro lado,
se $R(n_0, \dots, n_k)$ é falsa, então
$\Char{R}(n_0, \dots, n_k) = 0$. Isso significa que $\Th{Q}$ prova
\[
\lforall[y][(!A_{\Char{R}}(\num{n_0}, \dots, \num{n_k}, y) \lif y =
  \num{0})].
\]
Como $\Th{Q}$ prova $\eq/[\num{0}][\num{1}]$, $\Th{Q}$ prova
$\lnot !A_{\Char{R}}(\num{n_0}, \dots, \num{n_k}, \num{1})$ e,
portanto, prova $\lnot !A_R(\num{n_0}, \dots, \num{n_k})$.
\end{proof}

\begin{prob}
Mostre que, se $R$ é representável em~$\Th{Q}$, então~$\Char{R}$
também é.
\end{prob}

\end{document}
